% Consolidated 16-electrostatics-vacuum from 350eef03dcd4ca0e99eacbc61d97caa887d8c65a; original TeX fragments preserved.

% === BEGIN 16-electrostatics-vacuum-001-012.tex blob 736d5e713a7bdd3137a3cd1e9c18515c6901d5c1 ===
% Irodov Part III Electrodynamics §3.1, original 1--12.
% Source PDF pp.101--102 (printed pp.92--93), answers PDF pp.304--305.
% k=1/(4*pi*epsilon_0). Part III RESTARTS independent numbering.
\electroproblem{1}{92}
求两电子、两质子之间静电力与万有引力之比。
若两相同粒子的静电力和万有引力大小相等，
其电荷质量比 $|q|/m$ 应为多少？
\exsolution
库仑力与引力大小为
$F_e=e^2/(4\pi\varepsilon_0r^2)$、
$F_g=Gm^2/r^2$。
故比值与距离无关：
\[
\frac{F_e}{F_g}=\frac{e^2}{4\pi\varepsilon_0 Gm^2}.
\]
代入 $e=1.602\times10^{-19}\,\mathrm C$、
$m_e=9.109\times10^{-31}\,\mathrm{kg}$、
$m_p=1.673\times10^{-27}\,\mathrm{kg}$，
得到
\[
\boxed{(F_e/F_g)_{ee}\simeq4.2\times10^{42}},\quad
\boxed{(F_e/F_g)_{pp}\simeq1.2\times10^{36}}.
\]
一般相同粒子两力相等要求
$q^2/(4\pi\varepsilon_0r^2)=Gm^2/r^2$，
故
\[
\boxed{|q|/m=\sqrt{4\pi\varepsilon_0G}
\simeq8.6\times10^{-11}\,\mathrm{C/kg}}.
\]
\excheck \textbf{底本勘误：}
书末把质子比值印作 $10^{26}$，
与所列基本常数相差10个数量级；
应为约 $10^{36}$，两种力均按真空中的点粒子处理。

\electroproblem{2}{92}
两个铜球，质量各为 $1.0\,\mathrm g$，
球心相距 $1.0\,\mathrm m$。
若每球的电子总电荷与核内正电荷之差
为核电荷的 $1\%$，
求它们之间的电相互作用力数量级。
\exsolution
铜原子序数 $Z=29$、
摩尔质量 $M_{\rm Cu}\simeq63.5\,\mathrm{g/mol}$。
单球的原子数为
$N=(m/M_{\rm Cu})N_A$，
其核正电荷总量为 $ZNe$。
净电荷模
\[
|Q|=0.01\,ZNe
=0.01\cdot29\cdot\frac{1.0}{63.5}
(6.022\times10^{23})(1.602\times10^{-19})
\simeq4.4\times10^2\,\mathrm C.
\]
用库仑定律
\[
\boxed{F=\frac{Q^2}{4\pi\varepsilon_0r^2}
\simeq1.8\times10^{15}\,\mathrm N
\approx2\times10^{15}\,\mathrm N}.
\]
\excheck 这是概念性的巨大电荷差假设；
实际金属球不可能在普通环境下
稳定承载数百库仑的同号电荷。

\electroproblem{3}{92--93}
两相同小球，各质量 $m$、电荷 $q(t)$，
由长为 $\ell$ 的等长绝缘线悬挂于同一点。
两球距离 $x\ll\ell$。
若两球靠近速率满足
$-\dot x=v=a/x$，求电荷泄漏速率 $\dot q$。
\exsolution
在小角度准静态平衡下
$\tan\theta\simeq x/(2\ell)$，
两球间静电排斥力
$F_e=q^2/(4\pi\varepsilon_0x^2)$，
竖直平衡给线张力约为 $mg$，
因此
\[
\frac{q^2}{4\pi\varepsilon_0x^2}
=mg\frac{x}{2\ell}
\quad\Longrightarrow\quad
q=\sqrt{\frac{2\pi\varepsilon_0mg}{\ell}}\,
x^{3/2}.
\]
对时间求导并用 $\dot x=-a/x$：
\[
\boxed{\dot q=
-\frac{3a}{2}
\sqrt{\frac{2\pi\varepsilon_0mg}{\ell x}}}.
\]
泄漏量率为
$-\dot q=(3a/2)\sqrt{2\pi\varepsilon_0mg/(\ell x)}$。
\excheck 若原书抽取答案将 $x^{-1/2}$ 漏掉，
会导致量纲不等于 $\mathrm{C/s}$；
上式由逐步微分和单位检查确认。

\electroproblem{4}{93}
两个正点电荷 $q_1,q_2$，
位置矢量分别为 $\bm r_1,\bm r_2$。
求应放置在何处、大小为何的第三个负点电荷
$q_3$，使三个电荷各自受合力为零。
\exsolution
要使两正电荷的排斥力被吸引力抵消，
$q_3<0$ 必须处于两者连线之间。
设两正电荷间距 $L$，第三电荷距 $q_1$ 为 $x$。
第三电荷自身零受力要求
$q_1/x^2=q_2/(L-x)^2$，
得
\[
x=\frac{\sqrt{q_1}}{\sqrt{q_1}+\sqrt{q_2}}L,
\quad
\boxed{\bm r_3=
\frac{\sqrt{q_2}\,\bm r_1+\sqrt{q_1}\,\bm r_2}
{\sqrt{q_1}+\sqrt{q_2}}}.
\]
对 $q_1$ 的力平衡：
$q_1q_2/L^2=q_1|q_3|/x^2$，
故
\[
\boxed{q_3=-\frac{q_1q_2}
{(\sqrt{q_1}+\sqrt{q_2})^2}}.
\]
\excheck 对称特例 $q_1=q_2=q$，
$q_3=-q/4$ 且位于中点，
与库仑力的距离平方关系相符。

\electroproblem{5}{93}
半径 $r$ 的细导线圆环总电荷为 $q$。
若在圆环中心加入点电荷 $q_0$，
求导线的\emph{拉伸张力增量} $\Delta\mathcal T$。
\exsolution
环上每单位弧长的电荷
$\lambda=q/(2\pi r)$。
取中心角为 $\mathrm d\varphi$ 的微小圆环元，
其电荷
$\mathrm dq=\lambda r\,\mathrm d\varphi
=q\,\mathrm d\varphi/(2\pi)$。
中心电荷对该弧元施加径向力
$\mathrm dF=(1/(4\pi\varepsilon_0))
q_0\mathrm dq/r^2$。
两端张力增量的径向合力
$2\Delta\mathcal T\sin(\mathrm d\varphi/2)
\simeq\Delta\mathcal T\,\mathrm d\varphi$。
比较两式：
\[
\boxed{\Delta\mathcal T=
\frac{q q_0}{8\pi^2\varepsilon_0r^2}}.
\]
\excheck 相同符号电荷导致张力增加，
不同符号则为负增量；
这里求的是新增张力而非环自身静电排斥
原有的总张力。

\electroproblem{6}{93}
真空中有正点电荷 $q=50\,\mathrm{\mu C}$，
位于二维位置
$\bm r_0=(2\bm i+3\bm j)\,\mathrm m$。
求在
$\bm r=(8\bm i-5\bm j)\,\mathrm m$
处的电场矢量及模。
\exsolution
位移矢量
$\bm R=\bm r-\bm r_0
=(6\bm i-8\bm j)\,\mathrm m$，
其模 $R=10\,\mathrm m$。
点电荷电场
\[
\bm E=\frac{q}{4\pi\varepsilon_0}
\frac{\bm R}{R^3}.
\]
代入
$1/(4\pi\varepsilon_0)\simeq8.988\times10^9$
得
\[
\boxed{\bm E\simeq
(2.7\bm i-3.6\bm j)\,\mathrm{kV/m}},
\quad
\boxed{|\bm E|\simeq4.5\,\mathrm{kV/m}}.
\]
\excheck 负 $y$ 分量反映测量点位于电荷下方；
模满足 $\sqrt{2.7^2+3.6^2}=4.5$。

\electroproblem{7}{93}
图3.1所示的正方形两上顶点各放置 $+q$，
两下顶点各放置 $-q$，
正方形两条对角线长均为 $2\ell$。
求从正方形中心沿其法线距离 $x$ 处的
电场强度模。
\exsolution
取正方形位于 $z=0$，
上、下顶点纵坐标为
$\pm\ell/\sqrt2$，
各顶点到轴上观察点
$(0,0,x)$ 的距离均为
$\sqrt{\ell^2+x^2}$。
库仑场叠加时，
水平方向分量相互抵消，
轴向分量因总电荷为零而抵消，
只余从正电荷向负电荷方向的平面内分量。
单个顶点对该分量的贡献
为 $k q(\ell/\sqrt2)/(\ell^2+x^2)^{3/2}$，
四项相加：
\[
\boxed{E(x)=
\frac{q\ell}{\sqrt2\,\pi\varepsilon_0
(\ell^2+x^2)^{3/2}}}.
\]
\excheck 该场\emph{平行正方形所在平面}，
而不是垂直于平面；
$x\gg\ell$ 时呈偶极型 $x^{-3}$ 衰减。

\electroproblem{8}{93}
一条薄半圆弧半径 $R=20\,\mathrm{cm}$，
沿弧均匀带总电荷 $q=0.70\,\mathrm{nC}$。
求其曲率中心处电场强度模。
\exsolution
弧长 $\pi R$，线密度
$\lambda=q/(\pi R)$。
沿半圆以角度 $\varphi$ 参数化，
弧元
$\mathrm dq=\lambda R\,\mathrm d\varphi$，
电场大小微元为
$\mathrm dE=k\,\mathrm dq/R^2$。
平行弧直径的场分量两两抵消，
对称轴分量求和：
\[
E=\frac{k\lambda}{R}
\int_{-\pi/2}^{\pi/2}\cos\varphi\,\mathrm d\varphi
=\frac{2k\lambda}R
=\boxed{\frac{q}{2\pi^2\varepsilon_0R^2}}.
\]
代入得
\[
\boxed{E\simeq1.0\times10^2\,\mathrm{V/m}
=0.10\,\mathrm{kV/m}}.
\]
\excheck 完整均匀圆环的中心电场为零，
但半圆环沿其对称轴的场不为零。

\electroproblem{9}{93}
半径为 $r$ 的均匀带电细圆环总电荷为 $q$。
求圆环轴线上距中心 $\ell$ 处的电场模，
讨论 $\ell\gg r$，求最大值及其位置，
并说明 $E(\ell)$ 的大致图像。
\exsolution
每个环元到轴上点距离
$s=\sqrt{r^2+\ell^2}$，
径向分量由对称性相消，
仅保留轴向分量：
\[
E(\ell)=\frac{1}{4\pi\varepsilon_0}
\frac{|q\ell|}{(r^2+\ell^2)^{3/2}}.
\]
在 $\ell\gg r$ 时，
$E\simeq|q|/(4\pi\varepsilon_0\ell^2)$。
对 $\ell>0$ 求导：
\[
E'(\ell)\propto
\frac{r^2-2\ell^2}{(r^2+\ell^2)^{5/2}},
\]
故最大值出现在 $\ell=r/\sqrt2$，
\[
\boxed{E_{\max}=
\frac{|q|}{6\sqrt3\,\pi\varepsilon_0r^2}}.
\]
图像从 $E(0)=0$ 起先上升，
到 $\ell=r/\sqrt2$ 达峰值后按
$\ell^{-2}$ 衰减。
\excheck 环轴电场在中心恰为零；
若把环误当成点电荷，会在 $\ell=0$
得到错误的发散。

\electroproblem{10}{93}
一点电荷 $+q$ 处于均匀带电细圆环中心，
环半径为 $R$、总电荷为 $-q$。
求在轴线上 $x\gg R$ 处的电场模。
\exsolution
在 $x>0$ 处，
中心点电荷的轴向场
$E_+=kq/x^2$，
负圆环的轴向场
$E_-=-kq x/(x^2+R^2)^{3/2}$。
所以
\[
E_x=\frac{kq}{x^2}
\left[1-\left(1+\frac{R^2}{x^2}\right)^{-3/2}\right].
\]
当 $R/x\ll1$，
展开
$(1+u)^{-3/2}=1-\frac32u+O(u^2)$，
得到
\[
\boxed{E\simeq
\frac{3qR^2}{8\pi\varepsilon_0x^4}}.
\]
\excheck \textbf{原书书末系数争议：}
本版简答印
$3qR^2/(4\pi\varepsilon_0x^4)$，
是上述逐项叠加与二项式展开结果的两倍。
按照题面点电荷 $+q$ 与圆环总电荷 $-q$
的条件，本稿取 $3/(8\pi)$ 系数，
待独立对照原版勘误或其他版本。

\electroproblem{11}{93}
一个半径 $R$、总电荷为 $q$ 的均匀带电细圆环，
与一根从圆环中心沿轴向延伸至远方的
均匀带电细长直线同轴；
长直线从圆环中心开始，
线电荷密度为 $\lambda$。
求两者间的相互作用力。
\exsolution
从长直线起点起沿轴向取坐标 $z\ge0$。
线元电荷为 $\lambda\,\mathrm dz$，
到圆环任意环元的距离为
$\sqrt{R^2+z^2}$。
线元对圆环元的\emph{径向}场分量为
\[
\mathrm dE_R=
\frac1{4\pi\varepsilon_0}
\frac{\lambda R\,\mathrm dz}{(R^2+z^2)^{3/2}}.
\]
径向分量对整个均匀环对称相消；
更方便的办法是求\emph{轴向作用力}：
圆环轴线上电场为
$E_z(z)=kqz/(R^2+z^2)^{3/2}$，
长直线上受力为
\[
F=\lambda\int_0^\infty E_z(z)\,\mathrm dz
=kq\lambda
\int_0^\infty\frac{z\,\mathrm dz}{(R^2+z^2)^{3/2}}
=\boxed{\frac{q\lambda}{4\pi\varepsilon_0R}}.
\]
\excheck 作用力方向沿公共轴，
符号由 $q\lambda$ 决定；
使用牛顿第三定律，圆环受力大小相同方向相反。

\electroproblem{12}{93}
半径为 $R$ 的绝缘细圆环线电荷密度
$\lambda(\varphi)=\lambda_0\cos\varphi$。
求电场模：(a) 圆心；
(b) 环轴线上距圆心 $x$ 处；
并求 $x\gg R$ 的渐近表达式。
\exsolution
取 $\varphi=0$ 沿平面内 $+X$ 轴。
电荷元
$\mathrm dq=\lambda_0R\cos\varphi\,\mathrm d\varphi$，
$\int_0^{2\pi}\mathrm dq=0$。
观察点在轴线上时，场的轴向分量与总电荷成正比，
故为零；平面内 $X$ 分量为
\[
E_X=-k\frac{R}{(x^2+R^2)^{3/2}}
\int_0^{2\pi}\cos\varphi\,\mathrm dq
=-k\frac{\pi\lambda_0R^2}{(x^2+R^2)^{3/2}}.
\]
所以
\[
\boxed{E(0)=\frac{|\lambda_0|}{4\varepsilon_0R}},
\qquad
\boxed{E(x)=
\frac{|\lambda_0|R^2}{4\varepsilon_0(x^2+R^2)^{3/2}}}.
\]
定义环的偶极矩
$p=\int(R\cos\varphi)\,\mathrm dq
=\pi\lambda_0R^2$。
在 $x\gg R$ 时
\[
\boxed{E\simeq\frac{|p|}{4\pi\varepsilon_0x^3}}.
\]
\excheck 由于线密度正负交替，
该环\emph{净电荷为零}；
轴上场方向在环平面内，
表现为偶极场而非点电荷的 $x^{-2}$ 场。
% === END 16-electrostatics-vacuum-001-012.tex ===

% === BEGIN 16-electrostatics-vacuum-013-027.tex blob 1863ea21f67638fb68672f2619092a86ed59a7ce ===
% Irodov Part III §3.1, original 13--27. Source PDF pp.103--104,
% printed pp.94--95; answers PDF pp.305--306.
% Problems 17,18,25: textbook's printed coefficient conflicts with direct integration.
\electroproblem{13}{94}
一根长度为 $2a$ 的均匀带电细直杆总电荷为 $q$。
求距离杆中心为 $r$ 的点处电场模：
(a) 点在杆的中垂线上；
(b) 点在杆的延长线上（$r>a$）。
讨论 $r\gg a$。
\exsolution
设杆沿 $x$ 轴从 $-a$ 延伸至 $a$，
线电荷密度 $\lambda=q/(2a)$。
(a) 中垂线上点 $(0,r)$，横向分量抵消，
\[
E_\perp=k\lambda\int_{-a}^{a}
\frac{r\,\mathrm dx}{(r^2+x^2)^{3/2}}
=k\lambda\frac{2a}{r\sqrt{r^2+a^2}},
\]
因此
\[
\boxed{E_\perp=\frac{|q|}
{4\pi\varepsilon_0r\sqrt{r^2+a^2}}}.
\]
(b) 延长线上 $r>a$ 的点，
\[
E_\parallel=k\lambda\int_{-a}^{a}
\frac{\mathrm dx}{(r-x)^2}
=k\lambda\left(\frac1{r-a}-\frac1{r+a}\right),
\]
得到
\[
\boxed{E_\parallel=\frac{|q|}
{4\pi\varepsilon_0(r^2-a^2)}}.
\]
两种情况在 $r\gg a$ 时均有
$\boxed{E\simeq |q|/(4\pi\varepsilon_0r^2)}$。
\excheck 轴向表达式只在 $r>a$ 的杆外位置使用；
杆内理想线电荷的场需另外解释自场奇性。

\electroproblem{14}{94}
一根半无限长、均匀带电的直导线，
线密度为 $\lambda$，
求从其端点沿垂直方向距离为 $y$ 的点处
电场大小和方向。
\exsolution
取导线沿 $x\ge0$，观察点在 $(0,y)$。
电荷元 $\mathrm dq=\lambda\,\mathrm dx$，
位移向量为 $(-x,y)$。
因此
\[
E_x=-k\lambda\int_0^\infty
\frac{x\,\mathrm dx}{(x^2+y^2)^{3/2}}
=-\frac{k\lambda}{y},
\]
\[
E_y=k\lambda\int_0^\infty
\frac{y\,\mathrm dx}{(x^2+y^2)^{3/2}}
=\frac{k\lambda}{y}.
\]
合场大小
\[
\boxed{E=\frac{\sqrt2|\lambda|}
{4\pi\varepsilon_0y}}.
\]
\excheck 对正线电荷方向与直导线夹角为 $45^\circ$，
并从导线端点的上方指向远离导线的斜方向。

\electroproblem{15}{94}
一根很长、均匀线电荷密度 $\lambda$ 的细线
弯成图3.2的两种形状：
(a) 直角拐弯并由半径为 $R$ 的圆弧平滑连接；
(b) 半径为 $R$ 的对称 U 形弯曲。
求各图圆心 $O$ 的电场模；
直段长度远大于 $R$。
\exsolution
对每段细线用
\[
\mathrm d\bm E
=\frac{1}{4\pi\varepsilon_0}
\frac{\lambda\,\mathrm ds}{|\bm r_O-\bm r_s|^3}
(\bm r_O-\bm r_s)
\]
积分。长直段取半无限长极限。
(a) 直角弯的两直线段及圆弧产生的电场
按两个正交方向分解；
弧段的横、纵分量与直段部分抵消，
保留下来的两个分量大小均为 $k|\lambda|/R$。
故
\[
\boxed{E_a=\sqrt2\,\frac{k|\lambda|}{R}
=\frac{\sqrt2|\lambda|}{4\pi\varepsilon_0R}}.
\]
(b) U 形两侧直段的横向场分量因对称性相消，
其纵向和半圆弧的场积分大小相等、方向相反，
得到
\[
\boxed{E_b=0}.
\]
\excheck 结论与底本图3.2对应；
若把图(a)的两个圆弧中心误读为拐角点，
将得到错误的几何距离系数。

\electroproblem{16}{94}
半径为 $R$ 的球壳，表面电荷密度
$\sigma(\bm r)=\bm a\cdot\bm r$，
其中 $\bm a$ 为常矢量、
$\bm r$ 是从球心指向球面的矢量。
求球心处电场矢量。
\exsolution
令 $\bm r=R\bm n$，
$\mathrm dS=R^2\mathrm d\Omega$。
面元在球心产生的电场
\[
\mathrm d\bm E
=-k\frac{\sigma(\bm r)}{R^2}\bm n\,\mathrm dS
=-kR(\bm a\cdot\bm n)\bm n\,\mathrm d\Omega.
\]
利用球面角积分
$\int n_i n_j\,\mathrm d\Omega
=(4\pi/3)\delta_{ij}$，
得到
\[
\boxed{\bm E(0)=
-\frac{R}{3\varepsilon_0}\bm a}.
\]
\excheck 矢量方向与 $\bm a$ 相反，
大小和球壳半径 $R$ 成正比。

\electroproblem{17}{94}
半径为 $R$ 的球壳的面电荷分布为
$\sigma(\theta)=\sigma_0\cos\theta$，
$\theta$ 是对固定 $+z$ 轴的极角。
说明它可由两相反均匀带电实心球微小错位实现，
并求球壳内的电场矢量。
\exsolution
取两均匀体电荷密度为 $\pm\rho$ 的球
相对沿 $+z$ 轴错位 $\delta\ll R$。
两个球体表面相差的薄层在一阶近似中
形成面电荷密度
$\sigma(\theta)\simeq\rho\delta\cos\theta$，
故 $\rho\delta=\sigma_0$。
均匀带电实心球内部电场为
$\bm E=\rho(\bm r-\bm r_c)/(3\varepsilon_0)$。
正、负球的内部场叠加后，
$\bm r$ 项消掉，
\[
\bm E=
-\frac{\rho\delta}{3\varepsilon_0}\bm e_z
=\boxed{-\frac{\sigma_0}{3\varepsilon_0}\bm e_z}.
\]
\excheck \textbf{书末系数争议：}
本版简答印 $-\sigma_0\bm e_z/(2\varepsilon_0)$；
但直接球面积分
$-k\int\sigma_0\cos\theta\,\bm n\,\mathrm d\Omega$
和标准均匀极化球去极化场
均给系数 $1/3$，与题面所述球壳一致。
待核查书末是否混淆球壳与圆盘。

\electroproblem{18}{94}
半径为 $R$ 的实心球内，
体电荷密度
$\rho(\bm r)=\bm a\cdot\bm r$，
$\bm a$ 为常矢量。
求球心处电场。
\exsolution
取球坐标
$\bm r=r\bm n$，
体元
$\mathrm dV=r^2\mathrm dr\,\mathrm d\Omega$。
每一体元在球心的场为
\[
\mathrm d\bm E
=-k\frac{\rho(\bm r)}{r^2}\bm n\,\mathrm dV
=-k\,r(\bm a\cdot\bm n)\bm n\,
\mathrm dr\,\mathrm d\Omega.
\]
分离径向与角向积分：
\[
\bm E(0)
=-k\left(\int_0^Rr\,\mathrm dr\right)
\left(\int(\bm a\cdot\bm n)\bm n\,\mathrm d\Omega\right)
=-k\frac{R^2}{2}\frac{4\pi}{3}\bm a.
\]
故
\[
\boxed{\bm E(0)=
-\frac{R^2}{6\varepsilon_0}\bm a}.
\]
\excheck \textbf{书末系数争议：}
本版简答写 $-\bm aR^2/(2\varepsilon_0)$，
是上述体积分结果的3倍；
可由独立的对称张量平均
$\langle n_i n_j\rangle=\delta_{ij}/3$ 复核。

\electroproblem{19}{94}
一根半无限长均匀带电直线
沿半径为 $R$ 的圆盘轴线，
一端恰在圆盘中心，线密度为 $\lambda$。
求电场通过圆盘的通量，讨论法向选择。
\exsolution
在轴线上距离圆盘平面为 $z>0$ 的微元电荷
$\mathrm dq=\lambda\,\mathrm dz$，
圆盘张立体角大小为
$\Omega(z)=2\pi[1-z/\sqrt{z^2+R^2}]$。
对该点电荷，其电通量模为
$\mathrm d|\Phi|=(\mathrm dq/4\pi\varepsilon_0)\Omega(z)$。
于是
\[
|\Phi|=\frac{\lambda}{2\varepsilon_0}
\int_0^\infty
\left(1-\frac z{\sqrt{z^2+R^2}}\right)\mathrm dz.
\]
利用
$\int(1-z/\sqrt{z^2+R^2})\,\mathrm dz
=z-\sqrt{z^2+R^2}$，
在 $0,\infty$ 两端取值得 $R$，
因此
\[
\boxed{|\Phi|=\frac{|\lambda|R}{2\varepsilon_0}}.
\]
\excheck 圆盘是\emph{开曲面}，
通量符号取决于所选法向；
不能直接套闭曲面的高斯定理而忽视补面。

\electroproblem{20}{95}
图3.3中点电荷 $+q,-q$ 位于
半径为 $R$ 的圆盘两侧轴上，
分别距圆盘中心 $\ell$，
相距 $2\ell$。
求电场通过圆盘的通量。
\exsolution
每个点电荷在圆盘处张立体角
$\Omega=2\pi(1-\ell/\sqrt{\ell^2+R^2})$。
正电荷位于一侧，其场穿过圆盘；
负电荷位于另一侧，场线同向穿过，
两电荷贡献同号相加。
因此通量模为
\[
|\Phi|=\frac{2|q|}{4\pi\varepsilon_0}\Omega
=\boxed{\frac{|q|}{\varepsilon_0}
\left(1-\frac{\ell}{\sqrt{\ell^2+R^2}}\right)}.
\]
\excheck $R\to0$ 时通量趋零；
$R\to\infty$ 时通量模趋 $|q|/\varepsilon_0$；
正负号仍由圆盘法向定义决定。

\electroproblem{21}{95}
半径为 $R$ 的均匀体电荷密度
$\rho$ 的球体，
由距球心为 $r_0<R$ 的平面截出圆盘。
求电场穿过该截面的通量。
\exsolution
由均匀带电实心球的高斯定理，
球内位置矢量 $\bm r$ 的电场
$\bm E(\bm r)=\rho\bm r/(3\varepsilon_0)$。
取截面为 $z=r_0$，
截面内法向场分量处处相同：
$E_z=\rho r_0/(3\varepsilon_0)$。
圆盘半径
$b=\sqrt{R^2-r_0^2}$、
面积 $S=\pi(R^2-r_0^2)$。
因此
\[
\boxed{\Phi=
\frac{\pi\rho r_0(R^2-r_0^2)}{3\varepsilon_0}}.
\]
\excheck 若 $r_0=0$ 截面经过球心，
法向分量为零；若 $r_0=R$ 截面面积趋零，
通量同样趋零。

\electroproblem{22}{95}
两根很长的平行细线，
每根线电荷密度均为 $\lambda$，
相隔距离 $\ell$。
求位于两线中间的\emph{对称平面}内
电场强度的最大值。
\exsolution
令两线沿 $z$ 轴，
其横截面位置为
$x=\pm\ell/2,y=0$。
对称平面 $x=0$ 内，
两线产生的 $x$ 分量相消，
$y$ 分量相加：
\[
E_y(y)=
\frac{\lambda}{\pi\varepsilon_0}
\frac{y}{y^2+(\ell/2)^2}.
\]
令
$f(y)=|y|/[y^2+(\ell/2)^2]$，
对 $y>0$ 求导，
$f'(y)\propto(\ell/2)^2-y^2$。
极大出现在 $y=\ell/2$，故
\[
\boxed{E_{\max}=\frac{|\lambda|}{\pi\varepsilon_0\ell}}.
\]
\excheck 对称平面不是两线连线的平面；
沿二者中垂面 $x=0$，
恰在中点 $y=0$ 的场反而为零。

\electroproblem{23}{95}
无限长圆柱形薄壳上，
面电荷密度沿圆周方位角变化为
$\sigma(\varphi)=\sigma_0\cos\varphi$。
求圆柱轴线上的场强大小与方向。
\exsolution
取圆柱半径为 $R$，
观察点为轴线原点。
壳面元
$\mathrm dS=R\,\mathrm d\varphi\,\mathrm dz$，
电荷元
$\mathrm dq=\sigma_0\cos\varphi\,
R\,\mathrm d\varphi\,\mathrm dz$。
到轴线距离为
$\sqrt{R^2+z^2}$。
电场的横向 $\bm e_x$ 分量为
\[
E_x=-k\sigma_0 R^2
\int_0^{2\pi}\cos^2\varphi\,\mathrm d\varphi
\int_{-\infty}^{\infty}
\frac{\mathrm dz}{(R^2+z^2)^{3/2}}.
\]
两积分分别为 $\pi$ 和 $2/R^2$，
于是
\[
\boxed{\bm E=-\frac{\sigma_0}{2\varepsilon_0}\bm e_x}.
\]
\excheck 方向指向圆周上 $\varphi=\pi$ 的一侧，
场强与圆柱半径无关；
净线电荷密度沿圆周积分为零。

\electroproblem{24}{95}
某静电场只与 $x,y$ 有关：
\[
\bm E=a\,\frac{x\bm i+y\bm j}{x^2+y^2},
\]
其中 $a$ 常数。求通过球心位于原点、
半径为 $R$ 的球面的电通量。
\exsolution
在柱坐标中
$\bm E=(a/\varrho)\bm e_\varrho$，
是轴线附近具有奇性的径向场，
相当于等效无限线电荷场
$E_\varrho=\lambda/(2\pi\varepsilon_0\varrho)$。
匹配得 $\lambda=2\pi\varepsilon_0a$。
半径为 $R$ 的球面包围了轴线长为 $2R$
的一段有效线电荷，
$Q_{\rm enc}=\lambda(2R)$。
由高斯定律，
\[
\boxed{\Phi=Q_{\rm enc}/\varepsilon_0=4\pi aR}.
\]
\excheck 尽管 $\nabla\cdot\bm E=0$
在远离 $z$ 轴的点成立，
轴线处有分布意义的奇性；
不能由“处处散度零”误判球面通量为零。

\electroproblem{25}{95}
半径 $R$ 的绝缘球体电荷分布球对称：
$\rho(r)=\rho_0(1-r/R)$，$0\le r\le R$，
外部真空。求内外电场 $E(r)$、
最大值及其所处半径。
\exsolution
包围半径 $r\le R$ 内的总电荷为
\[
Q(r)=4\pi\int_0^r
\rho_0(1-s/R)s^2\,\mathrm ds
=4\pi\rho_0\left(\frac{r^3}3-\frac{r^4}{4R}\right).
\]
由高斯定律
$4\pi r^2E(r)=Q(r)/\varepsilon_0$，
\[
\boxed{E(r)=\frac{\rho_0r}{3\varepsilon_0}
\left(1-\frac{3r}{4R}\right)},\qquad r\le R.
\]
总电荷
$Q(R)=4\pi\rho_0R^3/12$，
所以外部
\[
\boxed{E(r)=\frac{\rho_0R^3}{12\varepsilon_0r^2}},
\qquad r\ge R.
\]
球内场求导
$E'(r)=\rho_0[1/3-r/(2R)]/\varepsilon_0$，
令其为零得
$r_{\max}=2R/3$，
\[
\boxed{E_{\max}=\frac{\rho_0R}{9\varepsilon_0}}.
\]
\excheck \textbf{本版书末第25(b)题系数疑误：}
简答写 $\rho_0R/(2\varepsilon_0)$，
代回其自身给出的 $E(r)$
在 $r=2R/3$ 处实际为
$\rho_0R/(9\varepsilon_0)$，后者由直接微分核实。

\electroproblem{26}{95}
半径为 $R$ 的球内部总电荷为 $q$，分布球对称；
球外介质从 $r=R$ 起另有空间电荷密度
$\rho(r)=\alpha/r$。
求让外部电场模与距离 $r$ 无关时，
球自身电荷 $q$ 的大小及该常场强。
\exsolution
由球对称的高斯定律，$r\ge R$ 时
\[
Q_{\rm enc}(r)
=q+4\pi\int_R^r(\alpha/s)s^2\,\mathrm ds
=q+2\pi\alpha(r^2-R^2).
\]
因此
\[
E(r)=\frac1{4\pi\varepsilon_0r^2}
[q-2\pi\alpha R^2+2\pi\alpha r^2].
\]
若 $E$ 与 $r$ 无关，
必须令 $1/r^2$ 的系数为零：
\[
\boxed{q=2\pi\alpha R^2},\qquad
\boxed{E=\frac{\alpha}{2\varepsilon_0}}.
\]
\excheck 电荷密度虽随 $1/r$ 减小，
球外包围的附加电荷仍按 $r^2$ 增长，
恰能维持常量场。

\electroproblem{27}{95}
空间中存在球对称电荷密度
$\rho(r)=\rho_0e^{-\alpha r^3}$，
$\rho_0,\alpha>0$。
求 $E(r)$，并研究
$\alpha r^3\ll1$ 与 $\alpha r^3\gg1$。
\exsolution
\textbf{题面图像复核：} 指数为
$-\alpha r^3$，PDF 文字抽取常丢失立方上标。
由高斯定律
\[
E(r)=\frac{1}{\varepsilon_0r^2}
\int_0^r\rho_0s^2e^{-\alpha s^3}\,\mathrm ds.
\]
令 $u=\alpha s^3$，则
$s^2\mathrm ds=\mathrm du/(3\alpha)$，
故
\[
\boxed{E(r)=\frac{\rho_0}
{3\varepsilon_0\alpha r^2}
(1-e^{-\alpha r^3})}.
\]
当 $\alpha r^3\ll1$，
$1-e^{-u}\simeq u$，所以
\[
\boxed{E(r)\simeq\frac{\rho_0r}{3\varepsilon_0}}.
\]
当 $\alpha r^3\gg1$，
指数项可略，得到
\[
\boxed{E(r)\simeq\frac{\rho_0}
{3\varepsilon_0\alpha r^2}}.
\]
\excheck 近球心呈均匀带电球内的线性规律；
远处呈有限总电荷产生的 $1/r^2$ 规律，
两个极限均正确。
% === END 16-electrostatics-vacuum-013-027.tex ===

% === BEGIN 16-electrostatics-vacuum-028-040.tex blob fa9cea68d33c71752bd2d44bb0e2b8c3c349873a ===
% Irodov Part III §3.1 original 28--40, PDF pp.104--106,
% printed pp.95--97, answer pages 306--307.
\electroproblem{28}{95--96}
一个均匀体电荷密度为 $\rho$ 的实心球内，
挖去一个球形空腔。
设空腔中心相对原球心的位移矢量为 $\bm a$。
若介电常数为1，求空腔内部电场。
\exsolution
用叠加原理：
“挖空”相当于
完整均匀带电球（密度 $\rho$）
与一颗在空腔位置、密度为 $-\rho$ 的球叠加。
在均匀带电球内部电场为
$\bm E=\rho\bm r/(3\varepsilon_0)$。
空腔内任一点的原球心位置为 $\bm r$，
相对空腔中心位置为 $\bm r-\bm a$，
故
\[
\bm E_{\rm cavity}
=\frac{\rho\bm r}{3\varepsilon_0}
-\frac{\rho(\bm r-\bm a)}{3\varepsilon_0}
=\boxed{\frac{\rho\bm a}{3\varepsilon_0}}.
\]
\excheck 空腔场为\emph{均匀场}，
与腔半径及腔内观察点位置无关；
$\bm a$ 按原球心指向腔心定义。

\electroproblem{29}{96}
均匀带电的无限长实心圆柱，
体电荷密度为 $\rho$。
内部挖有一条平行轴线的圆柱形空腔，
两轴相隔距离 $a$。
介电常数为1，求空腔内电场。
\exsolution
以截面位置矢量表示，
均匀带电无限长圆柱\emph{内部}电场为
$\bm E=\rho\bm r_\perp/(2\varepsilon_0)$。
挖空用密度为 $-\rho$ 的小圆柱叠加。
设位移 $\bm a$ 从原圆柱轴指向空腔轴，
则空腔内一点对小圆柱轴坐标为
$\bm r_\perp-\bm a$，
所以
\[
\boxed{\bm E_{\rm cavity}=
\frac{\rho\bm r_\perp}{2\varepsilon_0}
-\frac{\rho(\bm r_\perp-\bm a)}{2\varepsilon_0}
=\frac{\rho\bm a}{2\varepsilon_0}}.
\]
\excheck 空腔场与位置无关；
因圆柱体是二维横截面高斯对称，
系数为 $1/2$，
而球形空腔第28题系数为 $1/3$。

\electroproblem{30}{96}
两个同轴细圆环，半径相同为 $R$，
相距 $a$，分别带电 $+q,-q$。
求两个圆环中心之间的电势差。
\exsolution
细圆环总电荷为 $Q$ 时，
距环心轴向距离 $z$ 处电势为
$\varphi(z)=kQ/\sqrt{R^2+z^2}$。
设正环中心为 A、负环中心为 B，
则
\[
\varphi_A=k\left(\frac qR-\frac q{\sqrt{R^2+a^2}}\right),
\quad
\varphi_B=k\left(-\frac qR+\frac q{\sqrt{R^2+a^2}}\right).
\]
因此
\[
\boxed{\varphi_A-\varphi_B=
\frac q{2\pi\varepsilon_0R}
\left(1-\frac1{\sqrt{1+(a/R)^2}}\right)}.
\]
\excheck 当 $a\to0$，两个正负圆环重合，
电势差趋零；$a\to\infty$ 时跨中心电势差
趋 $q/(2\pi\varepsilon_0R)$。

\electroproblem{31}{96}
一根无限长均匀线电荷，线密度
$\lambda=0.40\,\mathrm{\mu C/m}$。
点2到导线的距离是点1的 $\eta=2.0$ 倍。
求两点电势差 $\varphi_1-\varphi_2$。
\exsolution
由高斯定律，径向电场
$E_r=\lambda/(2\pi\varepsilon_0r)$。
电势差为电场线积分的相反数：
\[
\varphi_1-\varphi_2=
\int_{r_1}^{r_2}E_r\,\mathrm dr
=\frac{\lambda}{2\pi\varepsilon_0}
\ln\frac{r_2}{r_1}
=\frac{\lambda}{2\pi\varepsilon_0}\ln\eta.
\]
代入题值：
\[
\boxed{\varphi_1-\varphi_2\simeq
5.0\times10^3\,\mathrm V=5.0\,\mathrm{kV}}.
\]
\excheck 无限线电荷\emph{绝对电势}
不能用无穷远为零来归一化，
但有限两点电势差有意义。

\electroproblem{32}{96}
半径 $R$ 的半球壳均匀面带电密度
$\sigma$。求球心处电势和电场模。
\exsolution
壳面上各点距球心都为 $R$，
半球面积 $2\pi R^2$，
故电势
\[
\varphi(0)=k\int_{\rm hemi}\frac{\sigma\,\mathrm dS}R
=k\frac{\sigma(2\pi R^2)}R
=\boxed{\frac{\sigma R}{2\varepsilon_0}}.
\]
设半球位于 $z>0$，
中心处场
$\mathrm d\bm E=-k\sigma\bm n\,\mathrm d\Omega$。
其轴向积分
$\int_{\rm hemi}\cos\theta\,\mathrm d\Omega
=2\pi\int_0^{\pi/2}\cos\theta\sin\theta\,\mathrm d\theta=\pi$，
因此
\[
\boxed{\bm E(0)=-\frac{\sigma}{4\varepsilon_0}\bm e_z},
\quad
\boxed{|\bm E|=\frac{|\sigma|}{4\varepsilon_0}}.
\]
\excheck 方向指向与带正电的半球相反一侧，
电势始终标量相加，无方向抵消。

\electroproblem{33}{96}
半径 $R$ 的无限薄均匀带电圆盘，
面电荷密度为 $\sigma$。
求沿圆盘轴线、距中心 $z>0$ 处的电势和场强；
分析 $z\to0^+$ 与 $z\gg R$。
\exsolution
把圆盘分解为半径 $s$、宽度 $\mathrm ds$
的同心圆环，
环电荷
$\mathrm dq=2\pi\sigma s\,\mathrm ds$。
轴上电势积分：
\[
\varphi(z)=k\int_0^R
\frac{2\pi\sigma s\,\mathrm ds}{\sqrt{z^2+s^2}}
=\boxed{\frac{\sigma}{2\varepsilon_0}
(\sqrt{z^2+R^2}-z)}.
\]
由 $\bm E=-\nabla\varphi$，
\[
\boxed{E_z(z)=\frac{\sigma}{2\varepsilon_0}
\left(1-\frac{z}{\sqrt{z^2+R^2}}\right)}.
\]
当 $z\to0^+$：
$\varphi\to\sigma R/(2\varepsilon_0)$、
$E_z\to\sigma/(2\varepsilon_0)$。
当 $z\gg R$，
$\sqrt{z^2+R^2}-z\simeq R^2/(2z)$；
设 $q=\pi R^2\sigma$，
\[
\boxed{\varphi\simeq\frac{q}{4\pi\varepsilon_0z},
\quad E_z\simeq\frac{q}{4\pi\varepsilon_0z^2}}.
\]
\excheck 在圆盘另一侧 $z<0$，
电势表达式须用 $|z|$，
电场法向方向随 $z$ 改变。

\electroproblem{34}{96}
半径为 $R$ 的薄圆盘均匀面带电密度
$\sigma$。求圆盘\emph{边缘}处电势
（无穷远电势取零）。
\exsolution
以圆盘边缘上的观察点作为极坐标原点。
取极轴指向圆盘中心，
一条与该方向成角 $\theta$ 的射线，
在圆盘内部的最长半径是
$s_{\max}=2R\cos\theta$，
$-\pi/2\le\theta\le\pi/2$。
从观察点到面电荷元的距离是 $s$，
面积元
$\mathrm dA=s\,\mathrm ds\,\mathrm d\theta$。
因此
\[
\varphi=k\sigma
\int_{-\pi/2}^{\pi/2}
\int_0^{2R\cos\theta}
\frac{s\,\mathrm ds\,\mathrm d\theta}{s}
=k\sigma(2R)\int_{-\pi/2}^{\pi/2}
\cos\theta\,\mathrm d\theta.
\]
积分后
\[
\boxed{\varphi_{\rm rim}=
4k\sigma R=\frac{\sigma R}{\pi\varepsilon_0}}.
\]
\excheck 圆盘边缘与中心不同；
中心电势是 $\sigma R/(2\varepsilon_0)$，
边缘电势较小，两者比值为 $2/\pi$。

\electroproblem{35}{96}
静电势为 $\varphi(\bm r)=\bm a\cdot\bm r$，
其中 $\bm a$ 为常矢量。
求电场强度。
\exsolution
写出笛卡尔分量
$\varphi=a_xx+a_yy+a_zz$，
其梯度
$\nabla\varphi=(a_x,a_y,a_z)=\bm a$。
由电势—电场关系
$\bm E=-\nabla\varphi$，
直接得到
\[
\boxed{\bm E=-\bm a}.
\]
\excheck 电场处处均匀，
方向与电势增长最快的方向相反。

\electroproblem{36}{96}
在 $(x,y)$ 平面中给定静电势：
(a) $\varphi=a(x^2-y^2)$；
(b) $\varphi=axy$。
求 $\bm E(x,y)$，并说明电场线的大致形状。
\exsolution
(a) 对电势求梯度：
\[
\boxed{\bm E_a=
-2ax\,\bm i+2ay\,\bm j}.
\]
场线微分方程
$\mathrm dy/\mathrm dx=E_y/E_x=-y/x$，
积分为
\[
\boxed{xy=\text{常数}}
\]
（坐标轴是分界线）。
(b) 同理
\[
\boxed{\bm E_b=-ay\,\bm i-ax\,\bm j}.
\]
场线
$\mathrm dy/\mathrm dx=x/y$，
积分为
\[
\boxed{y^2-x^2=\text{常数}}.
\]
两组场线都是双曲线族，
分别以坐标轴和斜线 $y=\pm x$ 为渐近线，
箭头方向由 $\bm E$ 的具体分量决定。
\excheck \textbf{本版书末第36(b)题符号疑误：}
简答将 $E_y$ 印作 $+ax$，
但由 $E_y=-\partial_y(axy)=-ax$ 必为负号。
若用书末符号，则
$\partial_xE_y-\partial_yE_x\ne0$，
也不符合静电场无旋条件。

\electroproblem{37}{96}
电势
$\varphi=a(x^2+y^2)+bz^2$。
求电场矢量及模，并讨论等势面
(a) $a>0,b>0$；
(b) $a>0,b<0$ 的形状。
\exsolution
直接对三坐标求负梯度：
\[
\boxed{\bm E=
-2a(x\bm i+y\bm j)-2bz\bm k},
\]
\[
\boxed{|\bm E|=
2\sqrt{a^2(x^2+y^2)+b^2z^2}}.
\]
等势面满足
$a(x^2+y^2)+bz^2=\varphi_*$。
(a) $a,b>0$ 且 $\varphi_*>0$，
是绕 $z$ 轴的旋转椭球面，
半轴 $\sqrt{\varphi_*/a}$、
$\sqrt{\varphi_*/b}$。
(b) $a>0,b<0$，
当 $\varphi_*>0$ 是单叶双曲面，
$\varphi_*=0$ 是圆锥面，
$\varphi_*<0$ 是双叶双曲面。
\excheck 等势面的分类依赖等势常数 $\varphi_*$
的符号，不可只依赖 $a,b$。

\electroproblem{38}{96}
半径 $R$ 的实心球内均匀分布总电荷 $q$，
球内外介电常数均为1。
求 (a) 球心电势；
(b) 距球心 $r<R$ 处电势。
\exsolution
由高斯定律，
球内场
$E(r)=qr/(4\pi\varepsilon_0R^3)$，
球外场
$E(r)=q/(4\pi\varepsilon_0r^2)$。
以无穷远电势为零，
球面电势
$\varphi(R)=q/(4\pi\varepsilon_0R)$。
球内
\[
\varphi(r)=\varphi(R)+\int_r^R E(s)\,\mathrm ds
=\frac{q}{4\pi\varepsilon_0R}
+\frac{q}{8\pi\varepsilon_0R^3}(R^2-r^2).
\]
因此
\[
\boxed{\varphi(r)=
\frac{q}{8\pi\varepsilon_0R}
\left(3-\frac{r^2}{R^2}\right)}.
\]
在球心 $r=0$：
\[
\boxed{\varphi(0)=\frac{3q}{8\pi\varepsilon_0R}}.
\]
\excheck 在球面 $r=R$ 与外部库仑电势连续，
且中心场 $-\varphi'(0)=0$。

\electroproblem{39}{97}
证明点偶极矩 $\bm p$ 在远场的电势为
$\varphi=k(\bm p\cdot\bm r)/r^3$，
并求电场模作为 $r,\theta$ 的函数，
其中 $\theta$ 是 $\bm p$ 与 $\bm r$ 夹角。
\exsolution
把偶极子视为两等量异号电荷
$+q,-q$，间距矢量为 $\bm d$，
令 $q|\bm d|\to p$ 而 $|\bm d|\to0$。
电势
\[
\varphi=kq\left(
\frac1{|\bm r-\bm d/2|}
-\frac1{|\bm r+\bm d/2|}\right)
\simeq kq\frac{\bm d\cdot\bm r}{r^3}
=\boxed{\frac{\bm p\cdot\bm r}
{4\pi\varepsilon_0r^3}}.
\]
球坐标中
$\varphi=kp\cos\theta/r^2$。
因此
$E_r=-\partial_r\varphi
=2kp\cos\theta/r^3$，
$E_\theta=-(1/r)\partial_\theta\varphi
=kp\sin\theta/r^3$。
合场模
\[
\boxed{E=\frac{|p|}{4\pi\varepsilon_0r^3}
\sqrt{1+3\cos^2\theta}}.
\]
\excheck 轴线上 $\theta=0$ 得
$E=2kp/r^3$；
赤道面 $\theta=\pi/2$ 得 $E=kp/r^3$，
轴上场为赤道场的两倍。

\electroproblem{40}{97}
点偶极矩 $\bm p=p\bm e_z$ 位于原点，
求电场沿 $z$ 轴的分量 $E_z$
及垂直 $z$ 轴方向的分量 $E_\perp$；
在何种角度电场垂直于 $\bm p$？
\exsolution
偶极场矢量形式
\[
\bm E=k\frac{3(\bm p\cdot\bm n)\bm n-\bm p}{r^3},
\qquad\bm n=\bm r/r.
\]
用
$\bm n=\sin\theta\,\bm e_\perp+
\cos\theta\,\bm e_z$ 分解：
\[
\boxed{E_z=\frac{p}{4\pi\varepsilon_0r^3}
(3\cos^2\theta-1)},
\]
\[
\boxed{E_\perp=\frac{3p}{4\pi\varepsilon_0r^3}
\sin\theta\cos\theta}.
\]
电场与 $\bm p$ 垂直等价于 $E_z=0$，
所以
\[
\cos^2\theta=\frac13,\qquad
\boxed{\theta=\arccos(1/\sqrt3)\simeq54.7^\circ
\quad\text{或}\quad125.3^\circ}.
\]
\excheck 本版书末第二角印约 $123.5^\circ$，
但应严格是
$180^\circ-54.7356^\circ=125.2644^\circ$；
可能存在数字排印错误。
% === END 16-electrostatics-vacuum-028-040.tex ===

% === BEGIN 16-electrostatics-vacuum-041-053.tex blob e36f092fa15046c7b54b3f04586a61ca3cbe666c ===
% Irodov Part III §3.1 original 41--53, PDF pp.106--107,
% printed pp.97--98; source figures 3.5--3.7 checked.
\electroproblem{41}{97}
点偶极矩 $\bm p$ 与均匀外电场 $\bm E_0$
方向相同。证明总电场有一张
以偶极子为球心的球形等势面，
求该球半径。
\exsolution
取 $z$ 轴沿 $\bm p$ 与 $\bm E_0$ 方向。
点偶极子电势
$\varphi_d=kp\cos\theta/r^2$，
均匀电场电势（选零点在原点）
$\varphi_0=-E_0r\cos\theta$。
总电势
\[
\varphi(r,\theta)=
\cos\theta\left(\frac{kp}{r^2}-E_0r\right).
\]
若括号为零，对\emph{全部}极角 $\theta$
有 $\varphi=0$，形成球形等势面。
由 $kp/R^2=E_0R$，
\[
\boxed{R=\left(\frac{p}{4\pi\varepsilon_0E_0}\right)^{1/3}}.
\]
\excheck 半径需取\emph{立方根}；
PDF 文字抽取容易把根号前的3写成乘数。

\electroproblem{42}{97}
两根相互平行、间距 $\ell$ 的长直细线，
线电荷密度分别为 $\lambda,-\lambda$。
求远离两线、距中点 $r\gg\ell$
且与连线夹角为 $\theta$ 处的电势和电场模。
\exsolution
在垂直两线的截面上，选两线中点为原点，
正负线位于 $\pm\bm \ell/2$。
二维线电荷的电势差叠加为
\[
\varphi=\frac{\lambda}{2\pi\varepsilon_0}
\ln\frac{r_-}{r_+}.
\]
在 $r\gg\ell$ 时，
$\ln(r_-/r_+)\simeq(\ell/r)\cos\theta$，
因此
\[
\boxed{\varphi(r,\theta)\simeq
\frac{\lambda\ell\cos\theta}{2\pi\varepsilon_0r}}.
\]
其极坐标场分量为
\[
E_r=-\partial_r\varphi=
\frac{\lambda\ell\cos\theta}{2\pi\varepsilon_0r^2},
\quad
E_\theta=-r^{-1}\partial_\theta\varphi=
\frac{\lambda\ell\sin\theta}{2\pi\varepsilon_0r^2}.
\]
故
\[
\boxed{E\simeq\frac{|\lambda|\ell}{2\pi\varepsilon_0r^2}}.
\]
\excheck 势依赖方向 $\cos\theta$，
但此二维偶极近场阶的\emph{场强模}
不依赖角度。

\electroproblem{43}{97}
两个同轴均匀带电细圆环，半径均为 $R$，
总电荷分别为 $+q,-q$，
中心相距 $\ell\ll R$。
按图3.6取中点为轴坐标 $x=0$，
求轴上电势 $\varphi(x)$ 和 $E_x(x)$
及其大致图形，并讨论 $|x|\gg R$。
\exsolution
令正环位于 $x=+\ell/2$，
负环位于 $x=-\ell/2$。
轴上电势精确表达式
\[
\varphi(x)=kq\left[
\frac1{\sqrt{R^2+(x-\ell/2)^2}}
-\frac1{\sqrt{R^2+(x+\ell/2)^2}}\right].
\]
对 $\ell/R\ll1$ 作一阶泰勒展开：
\[
\boxed{\varphi(x)\simeq
\frac{q\ell\,x}{4\pi\varepsilon_0
(R^2+x^2)^{3/2}}}.
\]
由 $E_x=-\mathrm d\varphi/\mathrm dx$ 得
\[
\boxed{E_x(x)\simeq
\frac{q\ell(2x^2-R^2)}
{4\pi\varepsilon_0(R^2+x^2)^{5/2}}}.
\]
$\varphi(x)$ 是奇函数，在 $x=0$ 过零，
在 $x=\pm R/\sqrt2$ 处有极值；
$E_x(x)$ 是偶函数，
中心为负，$|x|>R/\sqrt2$ 为正，
且远方趋于零。
当 $|x|\gg R$，
\[
\boxed{\varphi\sim\frac{q\ell\,\operatorname{sgn}x}
{4\pi\varepsilon_0x^2}},\qquad
\boxed{E_x\sim\frac{q\ell}{2\pi\varepsilon_0|x|^3}}.
\]
\begin{center}
\begin{tikzpicture}[scale=0.56,>=Stealth]
\draw[->](-3.4,0)--(3.4,0) node[right]{$x/R$};
\draw[->](0,-1.35)--(0,1.35) node[above]{$\varphi$};
\draw[thick,domain=-3:3,samples=60]
plot(\x,{2.6*\x/((1+\x*\x)^1.5)});
\begin{scope}[shift={(8,0)}]
\draw[->](-3.4,0)--(3.4,0) node[right]{$x/R$};
\draw[->](0,-1.35)--(0,1.35) node[above]{$E_x$};
\draw[thick,domain=-3:3,samples=60]
plot(\x,{1.8*(2*\x*\x-1)/((1+\x*\x)^2.5)});
\end{scope}
\end{tikzpicture}
\end{center}
\excheck 图线为无量纲形状示意。
远场是电偶极矩 $p=q\ell$ 沿环轴方向，
因此 $\varphi\propto |x|^{-2}$、
$E\propto|x|^{-3}$。

\electroproblem{44}{97}
两张无限大平面面电荷密度为
$+\sigma,-\sigma$，间距 $\ell\ll R$，
各有半径 $R$ 的同轴圆孔（图3.7）。
以孔轴中点为原点，沿轴取 $x$，
求电势 $\varphi(x)$ 与场分量 $E_x(x)$，
说明电势图像。
\exsolution
将“带圆孔平面”分解为
无限均匀带电平面，减去同号的带电圆盘。
两平面分离 $\ell\ll R$，
用薄层偶极展开并选
$\varphi(0)=0$，得到
\[
\boxed{\varphi(x)\simeq
\frac{\sigma\ell}{2\varepsilon_0}
\frac{x}{\sqrt{x^2+R^2}}}.
\]
对其求导：
\[
\boxed{E_x(x)=-\frac{\mathrm d\varphi}{\mathrm dx}
=-\frac{\sigma\ell R^2}
{2\varepsilon_0(x^2+R^2)^{3/2}}}.
\]
电势为单调增加的奇函数，
在 $x\to\pm\infty$ 时趋于
$\pm\sigma\ell/(2\varepsilon_0)$，
电场为关于原点对称的负向单峰。
\begin{center}
\begin{tikzpicture}[scale=0.67,>=Stealth]
\draw[->](-3.7,0)--(3.7,0) node[right]{$x/R$};
\draw[->](0,-1.3)--(0,1.3) node[above]{$\varphi$};
\draw[dashed](-3.3,1)--(3.3,1);
\draw[dashed](-3.3,-1)--(3.3,-1);
\draw[thick,domain=-3.3:3.3,samples=50]
plot(\x,{\x/sqrt(1+\x*\x)});
\end{tikzpicture}
\end{center}
\excheck 这里选择了图3.7中的面电荷符号；
若将正负平面互换，
$\varphi,E_x$ 整体反号，但模和函数形状不变。

\electroproblem{45}{98}
两张半径 $R$ 的平行圆盘相距
$\ell\ll R$，
面电荷密度分别为 $+\sigma,-\sigma$。
求沿轴线、离两圆盘距离
$x\gg\ell$ 处的电势和电场模，
再讨论 $x\gg R$。
\exsolution
一张均匀带电圆盘在轴上的电势为
$\varphi_{\rm disk}(z)=
[\sigma/(2\varepsilon_0)]
(\sqrt{R^2+z^2}-|z|)$。
分别对两个盘取正负号并用
$\ell\ll x$ 展开其轴向位移，
得到正盘外侧
\[
\boxed{\varphi_+(x)\simeq
\frac{\sigma\ell}{2\varepsilon_0}
\left(1-\frac{x}{\sqrt{x^2+R^2}}\right)},
\]
负盘外侧的电势取相反符号。
沿轴场强模由电势微分给出：
\[
\boxed{E(x)\simeq
\frac{|\sigma|\ell R^2}
{2\varepsilon_0(x^2+R^2)^{3/2}}}.
\]
远处 $x\gg R$，定义电偶极矩
$p=(\pi R^2\sigma)\ell$，于是
\[
\boxed{|\varphi|\simeq
\frac{|p|}{4\pi\varepsilon_0x^2}},\qquad
\boxed{E\simeq\frac{2|p|}{4\pi\varepsilon_0x^3}}.
\]
\excheck 圆盘尺寸仍远小于观察距离时才可
将有限平行板电容器视为点电偶极子。

\electroproblem{46}{98}
点偶极矩 $\bm p$ 位于距无限长均匀线电荷
$r$ 处，线电荷密度为 $\lambda$。
求作用力，当 $\bm p$ 分别：
(a) 沿直线；(b) 沿径向；
(c) 垂直于直线和径向。
\exsolution
线电荷外场为
$\bm E(\bm r_\perp)=
[\lambda/(2\pi\varepsilon_0)]\bm r_\perp/r^2$。
固定偶极矩的受力
$\bm F=(\bm p\cdot\nabla)\bm E$。
(a) 若沿线方向 $\bm p=p\bm e_z$，
场与 $z$ 无关，
\[
\boxed{\bm F_a=0}.
\]
(b) 若 $\bm p=p\bm e_r$，
径向求导
$\mathrm d(\lambda/(2\pi\varepsilon_0r))/\mathrm dr
=-\lambda/(2\pi\varepsilon_0r^2)$，
\[
\boxed{\bm F_b=
-\frac{\lambda p}{2\pi\varepsilon_0r^2}\bm e_r}.
\]
(c) 若 $\bm p=p\bm e_\varphi$，
在固定观察位置处沿切向微分
$\bm E$, 结果沿切向同向：
\[
\boxed{\bm F_c=
\frac{\lambda p}{2\pi\varepsilon_0r^2}\bm e_\varphi}.
\]
\excheck 受力是场梯度作用，
不是均匀场中的 $qE$；
若改变偶极矩的指向，受力也相应反向。

\electroproblem{47}{98}
两个水分子相距
$\ell=10\,\mathrm{nm}$，
各自电偶极矩大小
$p=0.62\times10^{-29}\,\mathrm{C\,m}$，
并沿两分子连线同向排列。
求两分子间相互作用力。
\exsolution
同轴同向的两点偶极子的相互作用势能为
\[
U(\ell)=
\frac{1}{4\pi\varepsilon_0\ell^3}
[\bm p_1\cdot\bm p_2
-3(\bm p_1\cdot\bm e_r)(\bm p_2\cdot\bm e_r)]
=-\frac{2p^2}{4\pi\varepsilon_0\ell^3}.
\]
相对连线方向的力
$F_r=-\mathrm dU/\mathrm d\ell$，
\[
F_r=-\frac{6p^2}{4\pi\varepsilon_0\ell^4}.
\]
所以力为吸引力，大小
\[
\boxed{|F|=
\frac{3p^2}{2\pi\varepsilon_0\ell^4}
\simeq2.1\times10^{-16}\,\mathrm N}.
\]
\excheck 假定两水分子可视为定向的点偶极子，
不是在自由转动分子的热平均取向下计算。

\electroproblem{48}{98}
电场给定为
$\bm E=a(y\bm i+x\bm j)$，
求电势 $\varphi(x,y)$。
\exsolution
由 $E_x=-\partial_x\varphi=-ay$，
对 $x$ 积分
$\varphi=-axy+f(y)$。
再用 $E_y=-\partial_y\varphi=-ax$，
得到 $f'(y)=0$，
因此
\[
\boxed{\varphi(x,y)=-axy+C}.
\]
\excheck $\partial_yE_x=\partial_xE_y=a$，
满足静电场无旋性，电势存在。

\electroproblem{49}{98}
电场
$\bm E=2axy\,\bm i+a(x^2-y^2)\bm j$，
求电势 $\varphi(x,y)$。
\exsolution
由
$\partial_x\varphi=-2axy$，
积分得
$\varphi=-ax^2y+f(y)$。
第二分量要求
$-\partial_y\varphi=a(x^2-y^2)$，
代入
$ax^2-f'(y)=ax^2-ay^2$，
所以 $f'(y)=ay^2$、
$f(y)=ay^3/3+C$。
最终
\[
\boxed{\varphi(x,y)=
ay\left(\frac{y^2}{3}-x^2\right)+C}.
\]
\excheck 两个场分量都由该电势的负梯度恢复，
可用 $\partial_yE_x=\partial_xE_y=2ax$
作独立检查。

\electroproblem{50}{98}
电场
\[
\bm E=ay\,\bm i+(ax+bz)\bm j+by\,\bm k,
\]
其中 $a,b$ 常数。求三维电势。
\exsolution
由 $\partial_x\varphi=-ay$，
先积分
$\varphi=-axy+f(y,z)$。
由 $E_y=-(\partial_y\varphi)
=ax+bz$，
得
$\partial_yf=-bz$，
所以
$f=-byz+g(z)$。
最后由
$E_z=-\partial_z\varphi=by-g'(z)$
得 $g'(z)=0$，
故
\[
\boxed{\varphi(x,y,z)=-y(ax+bz)+C}.
\]
\excheck 逐项负梯度给出原题全部三个分量，
表明输入场确为保守静电场。

\electroproblem{51}{98}
某区域电势只与 $x$ 有关：
$\varphi(x)=-ax^3+b$。
求该区域体电荷密度 $\rho(x)$。
\exsolution
泊松方程为
$\nabla^2\varphi=-\rho/\varepsilon_0$。
电势只依赖 $x$，所以
\[
\frac{\mathrm d^2\varphi}{\mathrm dx^2}
=\frac{\mathrm d}{\mathrm dx}(-3ax^2)=-6ax.
\]
由泊松方程
\[
\boxed{\rho(x)=-\varepsilon_0\nabla^2\varphi
=6\varepsilon_0ax}.
\]
\excheck 同时
$E_x=-\varphi_x=3ax^2$，
$\varepsilon_0\,\mathrm dE_x/\mathrm dx
=6\varepsilon_0ax$，与高斯定律一致。

\electroproblem{52}{98}
两块很大平行平板相距 $d$，
板间空间均匀填充体电荷密度 $\rho$，
板间电势差大小为 $|\Delta\varphi|$。
求使一块板附近电场强度为零的 $\rho$，
及另一板附近的场强大小。
\exsolution
沿板法线取 $x\in[0,d]$，
泊松方程等价于
$\mathrm dE_x/\mathrm dx=\rho/\varepsilon_0$。
若在 $x=0$ 处 $E_x(0)=0$，则
\[
E_x(x)=\frac{\rho x}{\varepsilon_0}.
\]
电势差
$\varphi(d)-\varphi(0)
=-\int_0^dE_x(x)\,\mathrm dx
=-\frac{\rho d^2}{2\varepsilon_0}$。
所以密度的\emph{大小}
\[
\boxed{|\rho|=
\frac{2\varepsilon_0|\Delta\varphi|}{d^2}},
\]
另一板邻近场强模为
\[
\boxed{|E(d)|=\frac{|\rho|d}{\varepsilon_0}
=\frac{2|\Delta\varphi|}{d}}.
\]
\excheck $\rho$ 的代数符号由
$\Delta\varphi$ 的端点定义确定；
本题若未指定电势差方向，
只能无歧义地报告其大小。

\electroproblem{53}{98}
带电实心球内部电势为
$\varphi(r)=ar^2+b$，
$a,b$ 为常数。
求球内空间电荷密度 $\rho(r)$。
\exsolution
径向球对称电势的拉普拉斯算符为
\[
\nabla^2\varphi=
\frac1{r^2}\frac{\mathrm d}{\mathrm dr}
\left(r^2\frac{\mathrm d\varphi}{\mathrm dr}\right).
\]
因为 $\varphi'(r)=2ar$，
\[
\nabla^2\varphi
=\frac1{r^2}\frac{\mathrm d(2ar^3)}{\mathrm dr}
=6a.
\]
由泊松方程
$\nabla^2\varphi=-\rho/\varepsilon_0$，
得到
\[
\boxed{\rho(r)=-6a\varepsilon_0=\text{常数}}.
\]
\excheck 常数电势偏移 $b$ 不产生电场或电荷；
体密度均匀，与球心正则性一致。
% === END 16-electrostatics-vacuum-041-053.tex ===
